
\documentclass{article}
\usepackage{colortbl}
\usepackage{makecell}
\usepackage{multirow}
\usepackage{supertabular}

\begin{document}

\newcounter{utterance}

\centering \large Interaction Transcript for game `dond', experiment `semi', episode 4 with Qwen3.6{-}35B{-}A3B{-}FP8{-}without{-}reasoning.
\vspace{24pt}

{ \footnotesize  \setcounter{utterance}{1}
\setlength{\tabcolsep}{0pt}
\begin{supertabular}{c@{$\;$}|p{.15\linewidth}@{}p{.15\linewidth}p{.15\linewidth}p{.15\linewidth}p{.15\linewidth}p{.15\linewidth}}
    \# & $\;$A & \multicolumn{4}{c}{Game Master} & $\;\:$B\\
    \hline

    \theutterance \stepcounter{utterance}  
    & & \multicolumn{4}{p{0.6\linewidth}}{
        \cellcolor[rgb]{0.9,0.9,0.9}{
            \makecell[{{p{\linewidth}}}]{
                \texttt{\tiny{[P1$\langle$GM]}}
                \texttt{أنت تلعب لعبة تفاوض تعاونية يتعين عليك فيها الاتفاق على كيفية تقسيم مجموعة من العناصر بينك وبين لاعب آخر.} \\
\\ 
\texttt{القواعد:} \\
\texttt{(a) يتم إعطاؤك أنت واللاعب الآخر مجموعة من العناصر. كما يُعطى كل منكما دالة قيمة سرية، تمثل مقدار تقديرك لكل نوع من الأشياء.} \\
\texttt{(b) تتبادل الرسائل مع اللاعب الآخر للاتفاق على من يحصل على أي عناصر. يمكن لكل منكما إرسال حد أقصى قدره 5 رسالة، أو الإنهاء مبكرًا عبر تقديم اقتراح سري في أي وقت.} \\
\texttt{(c) يُطلب من كل منكما تقديم اقتراح سري يحدد العناصر التي تريدها، منسقًا بين أقواس مربعة كما يلي: "{[}اقتراح: <رقم> <اسم الكائن>, <رقم> <اسم الكائن>, <...>{]}"} \\
\texttt{(d) إذا كانت اقتراحاتكما متكاملة، أي توجد عناصر كافية لتلبية كلا الاقتراحين، فسيُمنح كل لاعب درجة بناءً على مجموع القيم للعناصر التي حصل عليها. وإلا، فسيحصل كلاكما على صفر نقطة.} \\
\texttt{(e) هدفك هو تعظيم مجموع درجتك ودرجة اللاعب الآخر.} \\
\\ 
\texttt{لنبدأ.} \\
\\ 
\texttt{مجموعة العناصر المتاحة هي:} \\
\\ 
\texttt{1 كشتبان, 1 صابون, 4 طوابع} \\
\\ 
\texttt{دالة القيمة السرية الخاصة بك هي:} \\
\\ 
\texttt{كشتبان: 3, صابون: 7, طابع: 0} \\
\\ 
\texttt{مهم: رسائلك، ما لم تكن الاقتراح السري، تُنقل مباشرة إلى اللاعب الآخر، لذا لا تُضمّن أي رد على القواعد أو نصًا يعلن رسالتك. لتقديم اقتراح سري، استخدم التنسيق المشار إليه. لا تستخدم الأقواس المربعة عند التواصل مع اللاعب الآخر وإلا فسيُفسَّر ذلك على أنه اقتراحك السري.} \\
            }
        }
    }
    & \\ \\

    \theutterance \stepcounter{utterance}  
    & \multicolumn{4}{p{0.6\linewidth}}{
        \cellcolor[rgb]{0.9,0.9,0.9}{
            \makecell[{{p{\linewidth}}}]{
                \texttt{\tiny{[P1$\rangle$GM]}}
                \texttt{مرحبًا! يبدو أن لدينا فرصة جيدة للتوصل إلى اتفاق مفيد للطرفين.} \\
\\ 
\texttt{لديّ تقييم للعناصر كالتالي:} \\
\texttt{{-} الكشتبان: 3 نقاط لكل منهما (المجموع 6)} \\
\texttt{{-} الصابون: 7 نقاط} \\
\texttt{{-} الطوابع: 0 نقاط (لا قيمة لي لها)} \\
\\ 
\texttt{بما أن الطوابع لا قيمة لها بالنسبة لي، يمكنني التنازل عنها لك تمامًا دون أي تكلفة. الصابون له قيمة عالية بالنسبة لي (7 نقاط)، لذا سأحتاج إليه. الكشتبانان لديهما قيمة متوسطة (3 نقاط لكل منهما).} \\
\\ 
\texttt{أنا على استعداد للتفاوض. ما هي عناصرك المفضلة؟ وهل لديك تقييم مختلف للعناصر؟} \\
            }
        }
    }
    & & \\ \\

    \theutterance \stepcounter{utterance}  
    & & & \multicolumn{2}{p{0.3\linewidth}}{
        \cellcolor[rgb]{0.9,0.9,0.9}{
            \makecell[{{p{\linewidth}}}]{
                \texttt{\tiny{[GM$|$GM]}}
                \texttt{continue} \\
            }
        }
    }
    & & \\ \\

    \theutterance \stepcounter{utterance}  
    & & & \multicolumn{2}{p{0.3\linewidth}}{
        \cellcolor[rgb]{0.9,0.9,0.9}{
            \makecell[{{p{\linewidth}}}]{
                \texttt{\tiny{[GM$|$GM]}}
                \texttt{مرحبًا! يبدو أن لدينا فرصة جيدة للتوصل إلى اتفاق مفيد للطرفين.} \\
\\ 
\texttt{لديّ تقييم للعناصر كالتالي:} \\
\texttt{{-} الكشتبان: 3 نقاط لكل منهما (المجموع 6)} \\
\texttt{{-} الصابون: 7 نقاط} \\
\texttt{{-} الطوابع: 0 نقاط (لا قيمة لي لها)} \\
\\ 
\texttt{بما أن الطوابع لا قيمة لها بالنسبة لي، يمكنني التنازل عنها لك تمامًا دون أي تكلفة. الصابون له قيمة عالية بالنسبة لي (7 نقاط)، لذا سأحتاج إليه. الكشتبانان لديهما قيمة متوسطة (3 نقاط لكل منهما).} \\
\\ 
\texttt{أنا على استعداد للتفاوض. ما هي عناصرك المفضلة؟ وهل لديك تقييم مختلف للعناصر؟} \\
            }
        }
    }
    & & \\ \\

    \theutterance \stepcounter{utterance}  
    & & \multicolumn{4}{p{0.6\linewidth}}{
        \cellcolor[rgb]{0.9,0.9,0.9}{
            \makecell[{{p{\linewidth}}}]{
                \texttt{\tiny{[GM$\rangle$P2]}}
                \texttt{أنت تلعب لعبة تفاوض تعاونية يتعين عليك فيها الاتفاق على كيفية تقسيم مجموعة من العناصر بينك وبين لاعب آخر.} \\
\\ 
\texttt{القواعد:} \\
\texttt{(a) يتم إعطاؤك أنت واللاعب الآخر مجموعة من العناصر. كما يُعطى كل منكما دالة قيمة سرية، تمثل مقدار تقديرك لكل نوع من الأشياء.} \\
\texttt{(b) تتبادل الرسائل مع اللاعب الآخر للاتفاق على من يحصل على أي عناصر. يمكن لكل منكما إرسال حد أقصى قدره 5 رسالة، أو الإنهاء مبكرًا عبر تقديم اقتراح سري في أي وقت.} \\
\texttt{(c) يُطلب من كل منكما تقديم اقتراح سري يحدد العناصر التي تريدها، منسقًا بين أقواس مربعة كما يلي: "{[}اقتراح: <رقم> <اسم الكائن>, <رقم> <اسم الكائن>, <...>{]}"} \\
\texttt{(d) إذا كانت اقتراحاتكما متكاملة، أي توجد عناصر كافية لتلبية كلا الاقتراحين، فسيُمنح كل لاعب درجة بناءً على مجموع القيم للعناصر التي حصل عليها. وإلا، فسيحصل كلاكما على صفر نقطة.} \\
\texttt{(e) هدفك هو تعظيم مجموع درجتك ودرجة اللاعب الآخر.} \\
\\ 
\texttt{لنبدأ.} \\
\\ 
\texttt{مجموعة العناصر المتاحة هي:} \\
\\ 
\texttt{1 كشتبان, 1 صابون, 4 طوابع} \\
\\ 
\texttt{دالة القيمة السرية الخاصة بك هي:} \\
\\ 
\texttt{كشتبان: 2, صابون: 0, طابع: 2} \\
\\ 
\texttt{مهم: رسائلك، ما لم تكن الاقتراح السري، تُنقل مباشرة إلى اللاعب الآخر، لذا لا تُضمّن أي رد على القواعد أو نصًا يعلن رسالتك. لتقديم اقتراح سري، استخدم التنسيق المشار إليه. لا تستخدم الأقواس المربعة عند التواصل مع اللاعب الآخر وإلا فسيُفسَّر ذلك على أنه اقتراحك السري.} \\
\\ 
\texttt{مرحبًا! يبدو أن لدينا فرصة جيدة للتوصل إلى اتفاق مفيد للطرفين.} \\
\\ 
\texttt{لديّ تقييم للعناصر كالتالي:} \\
\texttt{{-} الكشتبان: 3 نقاط لكل منهما (المجموع 6)} \\
\texttt{{-} الصابون: 7 نقاط} \\
\texttt{{-} الطوابع: 0 نقاط (لا قيمة لي لها)} \\
\\ 
\texttt{بما أن الطوابع لا قيمة لها بالنسبة لي، يمكنني التنازل عنها لك تمامًا دون أي تكلفة. الصابون له قيمة عالية بالنسبة لي (7 نقاط)، لذا سأحتاج إليه. الكشتبانان لديهما قيمة متوسطة (3 نقاط لكل منهما).} \\
\\ 
\texttt{أنا على استعداد للتفاوض. ما هي عناصرك المفضلة؟ وهل لديك تقييم مختلف للعناصر؟} \\
            }
        }
    }
    & \\ \\

    \theutterance \stepcounter{utterance}  
    & & & \multicolumn{4}{p{0.6\linewidth}}{
        \cellcolor[rgb]{0.9,0.9,0.9}{
            \makecell[{{p{\linewidth}}}]{
                \texttt{\tiny{[GM$\langle$P2]}}
                \texttt{أهلاً بك! شكراً على عرضك الشفاف. يبدو أن لدينا مجالاً جيداً للتعاون لأن تفضيلاتنا تبدو متباينة إلى حد ما، مما يسهل الوصول إلى اتفاق يحقق قيمة عالية لكلا الطرفين.} \\
\\ 
\texttt{لأوضح لك وضعي بشكل دقيق:} \\
\texttt{{-} **الكشتبان:** لديّ قيمة قدرها **2** لكل منهما (المجموع 4).} \\
\texttt{{-} **الصابون:** لديّ قيمة **0** له تماماً (لا أهتم به).} \\
\texttt{{-} **الطوابع:** لديّ قيمة **2** لكل طابع (المجموع 8 لـ 4 طوابع).} \\
\\ 
\texttt{بناءً على ذلك، فإن أولويتي القصوى هي الحصول على **الطوابع**، تليها **الكشتبانان**، بينما **الصابون** غير مهم بالنسبة لي.} \\
\\ 
\texttt{بما أنك ذكرت أن الطوابع لا قيمة لها لديك، فهذا يعني أننا يمكننا منحها لك بالكامل دون أن تكلفك شيئاً، وفي المقابل ستحصل أنا على الكشتبانان أو الصابون حسب تفضيلاتك.} \\
\\ 
\texttt{لنقترح حلاً مبدئياً يحقق أقصى فائدة:} \\
\texttt{1. **أنا آخذ:** الكشتبانان (قيمتي: 4).} \\
\texttt{2. **أنت تأخذ:** الصابون (قيمتك: 7) وجميع الطوابع الأربعة (قيمتك: 0).} \\
\\ 
\texttt{في هذا السيناريو:} \\
\texttt{{-} درجتي: 4} \\
\texttt{{-} درجتك: 7} \\
\texttt{{-} المجموع الكلي: 11} \\
\\ 
\texttt{هل هذا الاقتراح مقبول لديك؟ إذا كنت تفضل الكشتبانان بدلاً من الصابون، يمكننا العكس، لكنني أعتقد أن هذا التوزيع يحقق توازناً جيداً حيث تحصل على أعلى عنصر قيمة لديك (الصابون) وأنا أحصل على ما أفضله (الكشتبانان).} \\
\\ 
\texttt{أنا على استعداد لتقديم اقتراحي السري الآن إذا وافقت على هذا التقسيم.} \\
            }
        }
    }
     \\ \\

    \theutterance \stepcounter{utterance}  
    & & & \multicolumn{2}{p{0.3\linewidth}}{
        \cellcolor[rgb]{0.9,0.9,0.9}{
            \makecell[{{p{\linewidth}}}]{
                \texttt{\tiny{[GM$|$GM]}}
                \texttt{continue} \\
            }
        }
    }
    & & \\ \\

    \theutterance \stepcounter{utterance}  
    & & & \multicolumn{2}{p{0.3\linewidth}}{
        \cellcolor[rgb]{0.9,0.9,0.9}{
            \makecell[{{p{\linewidth}}}]{
                \texttt{\tiny{[GM$|$GM]}}
                \texttt{أهلاً بك! شكراً على عرضك الشفاف. يبدو أن لدينا مجالاً جيداً للتعاون لأن تفضيلاتنا تبدو متباينة إلى حد ما، مما يسهل الوصول إلى اتفاق يحقق قيمة عالية لكلا الطرفين.} \\
\\ 
\texttt{لأوضح لك وضعي بشكل دقيق:} \\
\texttt{{-} **الكشتبان:** لديّ قيمة قدرها **2** لكل منهما (المجموع 4).} \\
\texttt{{-} **الصابون:** لديّ قيمة **0** له تماماً (لا أهتم به).} \\
\texttt{{-} **الطوابع:** لديّ قيمة **2** لكل طابع (المجموع 8 لـ 4 طوابع).} \\
\\ 
\texttt{بناءً على ذلك، فإن أولويتي القصوى هي الحصول على **الطوابع**، تليها **الكشتبانان**، بينما **الصابون** غير مهم بالنسبة لي.} \\
\\ 
\texttt{بما أنك ذكرت أن الطوابع لا قيمة لها لديك، فهذا يعني أننا يمكننا منحها لك بالكامل دون أن تكلفك شيئاً، وفي المقابل ستحصل أنا على الكشتبانان أو الصابون حسب تفضيلاتك.} \\
\\ 
\texttt{لنقترح حلاً مبدئياً يحقق أقصى فائدة:} \\
\texttt{1. **أنا آخذ:** الكشتبانان (قيمتي: 4).} \\
\texttt{2. **أنت تأخذ:** الصابون (قيمتك: 7) وجميع الطوابع الأربعة (قيمتك: 0).} \\
\\ 
\texttt{في هذا السيناريو:} \\
\texttt{{-} درجتي: 4} \\
\texttt{{-} درجتك: 7} \\
\texttt{{-} المجموع الكلي: 11} \\
\\ 
\texttt{هل هذا الاقتراح مقبول لديك؟ إذا كنت تفضل الكشتبانان بدلاً من الصابون، يمكننا العكس، لكنني أعتقد أن هذا التوزيع يحقق توازناً جيداً حيث تحصل على أعلى عنصر قيمة لديك (الصابون) وأنا أحصل على ما أفضله (الكشتبانان).} \\
\\ 
\texttt{أنا على استعداد لتقديم اقتراحي السري الآن إذا وافقت على هذا التقسيم.} \\
            }
        }
    }
    & & \\ \\

    \theutterance \stepcounter{utterance}  
    & & \multicolumn{4}{p{0.6\linewidth}}{
        \cellcolor[rgb]{0.9,0.9,0.9}{
            \makecell[{{p{\linewidth}}}]{
                \texttt{\tiny{[P1$\langle$GM]}}
                \texttt{أهلاً بك! شكراً على عرضك الشفاف. يبدو أن لدينا مجالاً جيداً للتعاون لأن تفضيلاتنا تبدو متباينة إلى حد ما، مما يسهل الوصول إلى اتفاق يحقق قيمة عالية لكلا الطرفين.} \\
\\ 
\texttt{لأوضح لك وضعي بشكل دقيق:} \\
\texttt{{-} **الكشتبان:** لديّ قيمة قدرها **2** لكل منهما (المجموع 4).} \\
\texttt{{-} **الصابون:** لديّ قيمة **0** له تماماً (لا أهتم به).} \\
\texttt{{-} **الطوابع:** لديّ قيمة **2** لكل طابع (المجموع 8 لـ 4 طوابع).} \\
\\ 
\texttt{بناءً على ذلك، فإن أولويتي القصوى هي الحصول على **الطوابع**، تليها **الكشتبانان**، بينما **الصابون** غير مهم بالنسبة لي.} \\
\\ 
\texttt{بما أنك ذكرت أن الطوابع لا قيمة لها لديك، فهذا يعني أننا يمكننا منحها لك بالكامل دون أن تكلفك شيئاً، وفي المقابل ستحصل أنا على الكشتبانان أو الصابون حسب تفضيلاتك.} \\
\\ 
\texttt{لنقترح حلاً مبدئياً يحقق أقصى فائدة:} \\
\texttt{1. **أنا آخذ:** الكشتبانان (قيمتي: 4).} \\
\texttt{2. **أنت تأخذ:** الصابون (قيمتك: 7) وجميع الطوابع الأربعة (قيمتك: 0).} \\
\\ 
\texttt{في هذا السيناريو:} \\
\texttt{{-} درجتي: 4} \\
\texttt{{-} درجتك: 7} \\
\texttt{{-} المجموع الكلي: 11} \\
\\ 
\texttt{هل هذا الاقتراح مقبول لديك؟ إذا كنت تفضل الكشتبانان بدلاً من الصابون، يمكننا العكس، لكنني أعتقد أن هذا التوزيع يحقق توازناً جيداً حيث تحصل على أعلى عنصر قيمة لديك (الصابون) وأنا أحصل على ما أفضله (الكشتبانان).} \\
\\ 
\texttt{أنا على استعداد لتقديم اقتراحي السري الآن إذا وافقت على هذا التقسيم.} \\
            }
        }
    }
    & \\ \\

    \theutterance \stepcounter{utterance}  
    & \multicolumn{4}{p{0.6\linewidth}}{
        \cellcolor[rgb]{0.9,0.9,0.9}{
            \makecell[{{p{\linewidth}}}]{
                \texttt{\tiny{[P1$\rangle$GM]}}
                \texttt{اقتراحك منطقي ويحقق توازنًا جيدًا. بما أن الطوابع لا قيمة لها لي، فأنا موافق على أخذها. الصابون له قيمة عالية لي (7)، والكشتبانان لهما قيمة متوسطة (3 لكل منهما).} \\
\\ 
\texttt{لنقم بتقسيم العناصر كالتالي:} \\
\texttt{{-} **أنت تأخذ:** الكشتبانان (قيمتك: 4) وجميع الطوابع الأربعة (قيمتك: 8).} \\
\texttt{{-} **أنا آخذ:** الصابون (قيمتي: 7).} \\
\\ 
\texttt{هذا التوزيع يحقق:} \\
\texttt{{-} درجتك: 4 + 8 = 12} \\
\texttt{{-} درجتي: 7} \\
\texttt{{-} المجموع الكلي: 19} \\
\\ 
\texttt{هذا أفضل من اقتراحك السابق (المجموع 11) لأننا نستخدم جميع العناصر ذات القيمة.} \\
\\ 
\texttt{سأقدم اقتراحي السري الآن:} \\
\\ 
\texttt{{[}اقتراح: 1 صابون{]}} \\
            }
        }
    }
    & & \\ \\

    \theutterance \stepcounter{utterance}  
    & & & \multicolumn{2}{p{0.3\linewidth}}{
        \cellcolor[rgb]{0.9,0.9,0.9}{
            \makecell[{{p{\linewidth}}}]{
                \texttt{\tiny{[GM$|$GM]}}
                \texttt{proposal} \\
            }
        }
    }
    & & \\ \\

    \theutterance \stepcounter{utterance}  
    & & & \multicolumn{2}{p{0.3\linewidth}}{
        \cellcolor[rgb]{0.9,0.9,0.9}{
            \makecell[{{p{\linewidth}}}]{
                \texttt{\tiny{[GM$|$GM]}}
                [0, 1, 0]
            }
        }
    }
    & & \\ \\

    \theutterance \stepcounter{utterance}  
    & & \multicolumn{4}{p{0.6\linewidth}}{
        \cellcolor[rgb]{0.9,0.9,0.9}{
            \makecell[{{p{\linewidth}}}]{
                \texttt{\tiny{[GM$\rangle$P2]}}
                \texttt{قدّم اللاعب الآخر اقتراحًا سريًا. حان الآن دورك لتقديم اقتراحك الخاص.} \\
            }
        }
    }
    & \\ \\

    \theutterance \stepcounter{utterance}  
    & & & \multicolumn{4}{p{0.6\linewidth}}{
        \cellcolor[rgb]{0.9,0.9,0.9}{
            \makecell[{{p{\linewidth}}}]{
                \texttt{\tiny{[GM$\langle$P2]}}
                \texttt{{[}اقتراح: 1 كشتبان, 2 كشتبان{]}} \\
            }
        }
    }
     \\ \\

    \theutterance \stepcounter{utterance}  
    & & & \multicolumn{2}{p{0.3\linewidth}}{
        \cellcolor[rgb]{0.9,0.9,0.9}{
            \makecell[{{p{\linewidth}}}]{
                \texttt{\tiny{[GM$|$GM]}}
                \texttt{proposal must include every item type only once} \\
            }
        }
    }
    & & \\ \\

\end{supertabular}
}

\end{document}
